\documentclass[aspectratio=169,11pt]{beamer}

% =========================================================
% Packages
% =========================================================
\usepackage{amsmath,amssymb,mathtools,bm}
\usepackage{booktabs}
\usepackage{array}
\usepackage{ragged2e}
\usepackage{microtype}
\usepackage{graphicx}
\usepackage{xcolor}

% =========================================================
% Theme
% =========================================================
\usetheme{Boadilla}
\usecolortheme{default}
\setbeamertemplate{navigation symbols}{}
\setbeamertemplate{footline}[frame number]

\definecolor{DeepBlue}{RGB}{28,66,120}
\definecolor{DarkGray}{RGB}{55,55,55}
\definecolor{SoftGray}{RGB}{242,242,242}
\definecolor{DarkRed}{RGB}{155,45,45}
\definecolor{DarkGreen}{RGB}{35,110,75}

\setbeamercolor{structure}{fg=DeepBlue}
\setbeamercolor{frametitle}{fg=DeepBlue}
\setbeamercolor{title}{fg=DeepBlue}
\setbeamercolor{block title}{bg=DeepBlue!12,fg=DeepBlue}
\setbeamercolor{block body}{bg=DeepBlue!3,fg=black}
\setbeamercolor{alerted text}{fg=DarkRed}

% =========================================================
% Layout tuning (compact vertical rhythm)
% =========================================================
\setbeamersize{text margin left=1.1em,text margin right=1.1em}

% Many slides contain several display equations; the default
% display skips are the single largest source of overfull frames.
\AtBeginDocument{%
  \setlength{\abovedisplayskip}{3.5pt plus 2pt minus 1.5pt}%
  \setlength{\belowdisplayskip}{3.5pt plus 2pt minus 1.5pt}%
  \setlength{\abovedisplayshortskip}{1.5pt plus 1pt minus 1pt}%
  \setlength{\belowdisplayshortskip}{2.5pt plus 1pt minus 1pt}%
  \setlength{\jot}{2pt}%
}

% Tighter blocks: trim the space around the box, not its inner padding
\addtobeamertemplate{block begin}{\vspace{-0.45em}}{}
\addtobeamertemplate{block end}{}{\vspace{-0.35em}}
\addtobeamertemplate{block alerted begin}{\vspace{-0.45em}}{}
\addtobeamertemplate{block alerted end}{}{\vspace{-0.35em}}

% Tighter lists
\setlength{\leftmargini}{1.5em}
\setbeamertemplate{itemize/enumerate body begin}{\vspace{-0.2em}}
\setbeamertemplate{itemize/enumerate body end}{\vspace{-0.2em}}

% Slightly tighter frame titles
\setbeamerfont{frametitle}{size=\large,series=\bfseries}

% Ragged-right table columns; narrow columns look much calmer when
% they are not hyphenated on nearly every line.
\newcommand{\NoHyph}{\hyphenpenalty=10000\exhyphenpenalty=10000\relax}
\newcolumntype{L}[1]{>{\RaggedRight\NoHyph\arraybackslash}p{#1}}
\newcolumntype{B}[1]{>{\RaggedRight\NoHyph\bfseries\arraybackslash}p{#1}}

% =========================================================
% Commands
% =========================================================
\newcommand{\E}{\mathbb{E}}
\newcommand{\Pp}{\mathbb{P}}
\newcommand{\ind}{\mathbf{1}}
\newcommand{\ATE}{\mathrm{ATE}}
\newcommand{\ATT}{\mathrm{ATT}}

\newcommand{\Question}{\textcolor{DarkRed}{\textbf{Question}}}
\newcommand{\Spec}{\textcolor{DeepBlue}{\textbf{Specification \& Estimand}}}
\newcommand{\Ident}{\textcolor{DarkGreen}{\textbf{Identification}}}
\newcommand{\Est}{\textcolor{violet}{\textbf{Estimation}}}
\newcommand{\Infer}{\textcolor{orange!80!black}{\textbf{Inference}}}



% =========================================================
% Title
% =========================================================
\title{Mostly Harmless Econometrics}
\subtitle{
\textbf{Session 2: What Does OLS Estimate?}
\\[0.55em]
Linear projection
$\rightarrow$ Weighted contrasts
$\rightarrow$ Weighted causal effects
$\rightarrow$ Estimand choice
}
\author{Kaiwen Luo}
\institute{Washington University in St.Louis}
\date{\today}

\begin{document}

% =========================================================
\begin{frame}
    \titlepage
\end{frame}


% ============================================================
% SECTION: What Does OLS Estimate?
% ============================================================

\section{What Does OLS Estimate?}

% ------------------------------------------------------------
\begin{frame}{What Does OLS Estimate? A Modern View}
\small

A useful way to think about OLS is:

\[
\boxed{
\text{Statistical estimand}
\quad
\xrightarrow{\text{additional assumptions}}
\quad
\text{causal estimand}
}
\]

\medskip

OLS does \alert{not} begin with a causal effect. It begins with a
population linear projection.

\medskip

As we impose additional structure, the same regression coefficient
may admit increasingly informative representations.

\begin{block}{Three questions}
For any OLS coefficient, ask:
\begin{enumerate}
    \item What is the coefficient \emph{algebraically}?
    \item Under what assumptions can it be written in terms of potential outcomes?
    \item Which treatment effects does it average, and with what weights?
\end{enumerate}
\end{block}

\end{frame}


% ------------------------------------------------------------
\begin{frame}{The Starting Point: No Causal Assumptions}
\small

Consider the population linear projection
\[
(\alpha^{LP},\beta^{LP})
=
\arg\min_{\alpha,\beta}
\mathbb{E}\left[(Y-\alpha-\beta D)^2\right]
\qquad\Longrightarrow\qquad
\boxed{
\beta^{LP}
=
\frac{\mathrm{Cov}(D,Y)}{\mathrm{Var}(D)}
}
\]
provided that \(\mathrm{Var}(D)>0\).

\medskip

\begin{columns}[T]
\begin{column}{0.42\textwidth}
This result requires:
\begin{itemize}
    \item no potential-outcome model;
    \item no exogeneity;
    \item no linear conditional expectation;
    \item no homogeneous treatment effect.
\end{itemize}
\end{column}
\begin{column}{0.54\textwidth}
\begin{block}{Interpretation}
\(\beta^{LP}\) is simply the slope of the best linear predictor of
\(Y\) using \(D\).
\end{block}
\end{column}
\end{columns}

\medskip

\[
\boxed{
\text{OLS is first a projection estimand, not a causal estimand.}
}
\]

\end{frame}


% ------------------------------------------------------------
\begin{frame}{The Assumption Ladder}
\small

It is useful to separate three layers.

\begin{block}{Layer 1: Pure algebra}
\[
\beta^{LP}
=
\frac{\mathrm{Cov}(D,Y)}
{\mathrm{Var}(D)}.
\]

This is always the population OLS estimand.
\end{block}

\begin{block}{Layer 2: Treatment support}
Knowing whether \(D\) is
\[
\text{binary},\qquad
\text{ordered discrete},\qquad
\text{or continuous}
\]
allows us to rewrite the same covariance ratio as an average of
different kinds of contrasts or slopes.
\end{block}

\begin{block}{Layer 3: Identification assumptions}
Assumptions such as
\[
\mathbb{E}[Y(d)\mid D,X]
=
\mathbb{E}[Y(d)\mid X]
\]
allow observed conditional contrasts to be interpreted causally.
\end{block}

\end{frame}


% ============================================================
% BINARY TREATMENT: NO COVARIATES
% ============================================================

\begin{frame}{Case 1: Binary Treatment, No Covariates}
\small

Suppose \(D\in\{0,1\}\). Then purely by algebra,

\[
\boxed{
\beta^{OLS}
=
\mathbb{E}[Y\mid D=1]
-
\mathbb{E}[Y\mid D=0].
}
\]

No causal assumption is needed for this equality.

\medskip

With potential outcomes
\(Y = D Y(1)+(1-D)Y(0)\),
we can write

\[
\begin{aligned}
\beta^{OLS}
&=
\mathbb{E}[Y(1)\mid D=1]
-
\mathbb{E}[Y(0)\mid D=0]
\\[2mm]
&=
\underbrace{
\mathbb{E}[Y(1)-Y(0)\mid D=1]
}_{ATT}
+
\underbrace{
\mathbb{E}[Y(0)\mid D=1]
-
\mathbb{E}[Y(0)\mid D=0]
}_{\text{selection bias}}.
\end{aligned}
\]

\end{frame}


% ------------------------------------------------------------
\begin{frame}{Binary Treatment: When Does OLS Become the ATE?}
\small

Suppose treatment is mean independent of potential outcomes:

\[
\mathbb{E}[Y(d)\mid D]
=
\mathbb{E}[Y(d)],
\qquad d\in\{0,1\}.
\]

Then

\[
\mathbb{E}[Y\mid D=1]
=
\mathbb{E}[Y(1)]
\qquad\text{and}\qquad
\mathbb{E}[Y\mid D=0]
=
\mathbb{E}[Y(0)],
\]

so that \(\beta^{OLS}=\mathbb{E}[Y(1)-Y(0)]=ATE\). Thus,

\[
\boxed{
\frac{\mathrm{Cov}(D,Y)}{\mathrm{Var}(D)}
\overset{D\in\{0,1\}}{=}
\mathbb{E}[Y\mid 1]-\mathbb{E}[Y\mid 0]
\overset{\text{exogeneity}}{=}
ATE.
}
\]

\begin{block}{Key point}
The first equality is algebra.
The second equality is identification.
\end{block}

\end{frame}


% ============================================================
% COVARIATES: GENERAL SETUP
% ============================================================

\begin{frame}{Adding Covariates: Two Different Objects}
\small

Consider the familiar additive regression
\(Y=\alpha+\beta D+X'\gamma+u\).
By Frisch--Waugh--Lovell,

\[
\boxed{
\beta^{OLS}
=
\frac{\mathbb{E}[\widetilde D Y]}
{\mathbb{E}[\widetilde D^2]},
\qquad
\widetilde D
=
D-L(D\mid 1,X),
}
\]

where \(L(D\mid 1,X)\) is the \emph{linear projection} of \(D\) on \(X\).

\medskip

This is generally different from residualizing by the true conditional mean,
\(D-\mathbb{E}[D\mid X]\).

\begin{alertblock}{Important distinction}
In general,
\[
L(D\mid 1,X)
\neq
\mathbb{E}[D\mid X].
\]

Therefore, ordinary additive OLS with a finite-dimensional set of
controls does not automatically inherit the clean causal weighting
formulas that arise under saturated or flexible adjustment.
\end{alertblock}

\end{frame}


% ------------------------------------------------------------
\begin{frame}{A Clean Benchmark: Flexible Partialling Out}
\small

Consider instead the partially linear projection

\[
(\beta^*,g^*)
=
\arg\min_{\beta,g}
\mathbb{E}
\left[
\left(Y-\beta D-g(X)\right)^2
\right],
\]

where \(g(X)\) is unrestricted.
Define \(\mu_D(X)=\mathbb{E}[D\mid X]\). Then

\[
\boxed{
\beta^*
=
\frac{
\mathbb{E}
\left[
(D-\mu_D(X))Y
\right]
}{
\mathbb{E}
\left[
(D-\mu_D(X))^2
\right]
}.
}
\]

Using iterated expectations,

\[
\boxed{
\beta^*
=
\frac{
\mathbb{E}
\left[
\mathrm{Cov}(D,Y\mid X)
\right]
}{
\mathbb{E}
\left[
\mathrm{Var}(D\mid X)
\right]
}.
}
\]

This provides a useful benchmark for all treatment types.

\end{frame}


% ============================================================
% BINARY TREATMENT + COVARIATES
% ============================================================

\begin{frame}{Case 2: Binary Treatment with Flexible Covariate Adjustment}
\small

Let

\[
D\in\{0,1\},
\qquad
e(X)=\Pr(D=1\mid X).
\]

Then

\[
\mathrm{Var}(D\mid X)
=
e(X)\{1-e(X)\}.
\]

Also,

\[
\mathrm{Cov}(D,Y\mid X)
=
e(X)\{1-e(X)\}
\left[
\mathbb{E}(Y\mid D=1,X)
-
\mathbb{E}(Y\mid D=0,X)
\right].
\]

Hence

\[
\boxed{
\beta^*
=
\frac{
\mathbb{E}
\left[
e(X)\{1-e(X)\}
\Delta_Y(X)
\right]
}{
\mathbb{E}
\left[
e(X)\{1-e(X)\}
\right]
},
}
\]

where

\[
\Delta_Y(X)
=
\mathbb{E}[Y\mid D=1,X]
-
\mathbb{E}[Y\mid D=0,X].
\]

\end{frame}


% ------------------------------------------------------------
\begin{frame}{Binary Treatment + Conditional Ignorability}
\small

Suppose conditional mean independence holds together with overlap:

\[
\mathbb{E}[Y(d)\mid D,X]
=
\mathbb{E}[Y(d)\mid X],
\quad d\in\{0,1\},
\qquad\qquad
0<e(X)<1.
\]

Define the conditional average treatment effect

\[
\tau(X)
=
\mathbb{E}[Y(1)-Y(0)\mid X].
\]

Then \(\Delta_Y(X)=\tau(X)\), and therefore

\[
\boxed{
\beta^*
=
\frac{
\mathbb{E}
\left[
e(X)\{1-e(X)\}\tau(X)
\right]
}{
\mathbb{E}
\left[
e(X)\{1-e(X)\}
\right]
}.
}
\]

\begin{block}{Interpretation}
The coefficient is an
\alert{overlap-weighted average treatment effect},
not generally the population ATE.
\end{block}

\end{frame}


% ------------------------------------------------------------
\begin{frame}{Binary Treatment: Why the OLS Weight Matters}
\small

The weight is proportional to \(e(X)\{1-e(X)\}\).
It is largest when \(e(X)=\tfrac{1}{2}\), and approaches zero when
\(e(X)\rightarrow 0\) or \(e(X)\rightarrow 1\).

\medskip

Thus the estimand emphasizes covariate cells in which both treated
and untreated observations are common.

\medskip

In general,

\[
\boxed{
\beta^*
\neq
ATE
=
\mathbb{E}[\tau(X)].
}
\]

Equality holds under additional conditions, for example:

\begin{itemize}
    \item homogeneous treatment effects, \(\tau(X)=\tau\);
    \item randomized treatment with constant propensity, \(e(X)=p\);
    \item or other special restrictions that make the OLS weights
    unrelated to treatment-effect heterogeneity.
\end{itemize}

\end{frame}


% ============================================================
% ORDERED DISCRETE TREATMENT
% ============================================================

\begin{frame}{Case 3: Ordered Discrete Treatment, No Covariates}
\small

Suppose

\[
D\in\{d_0,d_1,\ldots,d_J\},
\qquad
d_0<d_1<\cdots<d_J,
\]

and define the conditional mean
\(m(d)=\mathbb{E}[Y\mid D=d]\).

\medskip

Since \(\mathbb{E}[Y-m(D)\mid D]=0\), the residual \(Y-m(D)\) is
uncorrelated with \(D\), and we obtain

\[
\boxed{
\beta^{OLS}
=
\frac{
\mathrm{Cov}(D,m(D))
}{
\mathrm{Var}(D)
}.
}
\]

\begin{block}{Implication}
OLS only uses the conditional mean function \(m(\cdot)\).
The question becomes: how does the covariance ratio summarize the
shape of \(m(\cdot)\) across the support of \(D\)?
\end{block}

\end{frame}


% ------------------------------------------------------------
\begin{frame}{Case 3: Ordered Discrete Treatment as Weighted Slopes}
\small

Define the adjacent slope

\[
s_j
=
\frac{
m(d_j)-m(d_{j-1})
}{
d_j-d_{j-1}
},
\qquad j=1,\ldots,J.
\]

Then

\[
\boxed{
\beta^{OLS}
=
\sum_{j=1}^{J} w_j s_j,
}
\qquad
w_j
=
\frac{
(d_j-d_{j-1})
\,\mathrm{Cov}
\left(
D,\mathbf{1}\{D\geq d_j\}
\right)
}{
\mathrm{Var}(D)
}.
\]

Moreover,

\[
w_j\geq 0,
\qquad
\sum_{j=1}^{J}w_j=1.
\]

\begin{block}{Interpretation}
OLS is a convex combination of adjacent slopes of the conditional
mean function. Margins near the middle of the distribution of \(D\)
typically receive the largest weights.
\end{block}

\end{frame}


% ------------------------------------------------------------
\begin{frame}{Ordered Discrete Treatment: Causal Interpretation}
\small

Suppose
\(\mathbb{E}[Y(d)\mid D]=\mathbb{E}[Y(d)]\) for every \(d\).
Then \(m(d)=\mathbb{E}[Y(d)]\). Therefore,

\[
\boxed{
\beta^{OLS}
=
\sum_{j=1}^{J}
w_j
\frac{
\mathbb{E}[Y(d_j)-Y(d_{j-1})]
}{
d_j-d_{j-1}
}.
}
\]

\begin{block}{Interpretation}
OLS is a weighted average of causal effects associated with
\emph{different treatment margins}:
\[
d_0\rightarrow d_1,\quad
d_1\rightarrow d_2,\quad
\ldots,\quad
d_{J-1}\rightarrow d_J.
\]
\end{block}

Unlike the binary case, there is no longer a unique treatment margin.

\medskip

\scriptsize
Related result: Yitzhaki (1996), \emph{Journal of Business \& Economic Statistics}.

\end{frame}


% ------------------------------------------------------------
\begin{frame}{Ordered Discrete Treatment with Covariates}
\small

Under flexible partialling out,

\[
\beta^*
=
\frac{
\mathbb{E}
[
\mathrm{Var}(D\mid X)\beta(X)
]
}{
\mathbb{E}
[
\mathrm{Var}(D\mid X)
]
},
\qquad
\beta(X)
=
\frac{
\mathrm{Cov}(D,Y\mid X)
}{
\mathrm{Var}(D\mid X)
}.
\]

Within each covariate cell, with \(m(d,X)=\mathbb{E}[Y\mid D=d,X]\),

\[
\beta(X)
=
\sum_{j=1}^{J}
w_j(X)
\frac{
m(d_j,X)-m(d_{j-1},X)
}{
d_j-d_{j-1}
}.
\]

Thus OLS contains \alert{two layers of weighting}:

\[
\boxed{
\underbrace{\mathrm{Var}(D\mid X)}_{\text{across covariate cells}}
\times
\underbrace{w_j(X)}_{\text{across treatment margins}}.
}
\]

\end{frame}


% ------------------------------------------------------------
\begin{frame}{Ordered Discrete Treatment + Conditional Ignorability}
\small

Suppose
\(\mathbb{E}[Y(d)\mid D,X]=\mathbb{E}[Y(d)\mid X]\).
Then \(m(d,X)=\mathbb{E}[Y(d)\mid X]\).

\medskip

Define the conditional causal slope

\[
\tau_j(X)
=
\frac{
\mathbb{E}[Y(d_j)-Y(d_{j-1})\mid X]
}{
d_j-d_{j-1}
}.
\]

Then

\[
\boxed{
\beta^*
=
\frac{
\mathbb{E}
\left[
\mathrm{Var}(D\mid X)
\sum_{j=1}^{J}
w_j(X)\tau_j(X)
\right]
}{
\mathbb{E}
[
\mathrm{Var}(D\mid X)
]
}.
}
\]

\begin{block}{Interpretation}
A variance-weighted average across covariate cells,
within which OLS further averages across treatment margins.
\end{block}

\end{frame}


% ------------------------------------------------------------
\begin{frame}{A Caveat: Unordered Discrete Treatments}
\small

Suppose instead that treatment categories have no natural ordering:
\(D\in\{\text{A},\text{B},\text{C}\}\).

\medskip

Assigning arbitrary numerical codes such as
\(A=1,\ B=2,\ C=3\)
and estimating a single regression slope is generally not meaningful.
Changing the coding changes the estimand.

\medskip

The natural approach is to use treatment indicators:

\[
Y
=
\alpha
+
\beta_B\mathbf{1}\{D=B\}
+
\beta_C\mathbf{1}\{D=C\}
+
u.
\]

Without covariates,

\[
\beta_B
=
\mathbb{E}[Y\mid D=B]
-
\mathbb{E}[Y\mid D=A],
\]

and similarly for \(\beta_C\).

\begin{block}{Conclusion}
A scalar ``OLS slope'' is intrinsically meaningful for an
\emph{ordered scalar treatment}, but not for arbitrarily coded
unordered categories.
\end{block}

\end{frame}


% ============================================================
% CONTINUOUS TREATMENT
% ============================================================

\begin{frame}{Case 4: Continuous Treatment, No Covariates}
\small

Suppose \(D\) is continuously distributed on
\(\mathcal{D}=[a,b]\), and define
\(m(d)=\mathbb{E}[Y\mid D=d]\). Again,

\[
\beta^{OLS}
=
\frac{
\mathrm{Cov}(D,m(D))
}{
\mathrm{Var}(D)
}.
\]

If \(m(d)\) is differentiable, then

\[
\boxed{
\beta^{OLS}
=
\int_a^b
\omega(d)m'(d)\,dd,
}
\qquad
\boxed{
\omega(d)
=
\frac{
\mathrm{Cov}
\left(
D,\mathbf{1}\{D\geq d\}
\right)
}{
\mathrm{Var}(D)
}.
}
\]

The weights satisfy

\[
\omega(d)\geq 0,
\qquad
\int_a^b\omega(d)\,dd=1.
\]

\begin{block}{Interpretation}
The continuous analogue of the ordered discrete case: adjacent
slopes become local derivatives \(m'(d)\).
\end{block}

\end{frame}


% ------------------------------------------------------------
\begin{frame}{Continuous Treatment: What Does the Slope Average?}
\small

The OLS coefficient is therefore a weighted average derivative:

\[
\boxed{
\beta^{OLS}
=
\int
\omega(d)
\frac{\partial}{\partial d}
\mathbb{E}[Y\mid D=d]
\,dd.
}
\]

This should be distinguished from the ordinary average derivative

\[
ADE
=
\mathbb{E}[m'(D)]
=
\int
m'(d)f_D(d)\,dd.
\]

In general,

\[
\boxed{
\omega(d)\neq f_D(d)
}
\qquad\Longrightarrow\qquad
\boxed{
\beta^{OLS}\neq ADE.
}
\]

\begin{block}{Key lesson}
With a nonlinear conditional mean, a linear regression coefficient
is one particular weighted summary of local slopes.
\end{block}

\scriptsize
Related result: Yitzhaki (1996).

\end{frame}


% ------------------------------------------------------------
\begin{frame}{Continuous Treatment: Causal Interpretation}
\small

Suppose
\(\mathbb{E}[Y(d)\mid D]=\mathbb{E}[Y(d)]\) for every \(d\).
Then \(m(d)=\mathbb{E}[Y(d)]\). Hence

\[
\boxed{
\beta^{OLS}
=
\int
\omega(d)
\frac{\partial}{\partial d}
\mathbb{E}[Y(d)]
\,dd.
}
\]

Thus the coefficient can be interpreted as a weighted average of
marginal causal responses at different treatment levels.

\medskip

But this does \alert{not} imply that OLS estimates a uniquely
privileged ``average causal effect.''

\begin{block}{Important}
For continuous treatments, different scientifically meaningful
summaries include:
\[
\text{dose-response curves},\quad
\text{shift effects},\quad
\text{average derivative effects},\quad
\text{weighted derivative effects}.
\]
There is no unique scalar analogue of the binary-treatment ATE.
\end{block}

\end{frame}


% ------------------------------------------------------------
\begin{frame}{Case 5: Continuous Treatment with Covariates}
\small

Under flexible partialling out,

\[
\beta^*
=
\frac{
\mathbb{E}
[
\mathrm{Var}(D\mid X)\beta(X)
]
}{
\mathbb{E}
[
\mathrm{Var}(D\mid X)
]
},
\qquad
\beta(X)
=
\frac{
\mathrm{Cov}(D,Y\mid X)
}{
\mathrm{Var}(D\mid X)
}.
\]

For a differentiable conditional mean
\(m(d,x)=\mathbb{E}[Y\mid D=d,X=x]\),
we can write

\[
\boxed{
\beta(x)
=
\int_{\mathcal{D}(x)}
\omega(d\mid x)
\frac{\partial m(d,x)}{\partial d}
\,dd,
}
\]

with

\[
\omega(d\mid x)
=
\frac{
\mathrm{Cov}
\left(
D,\mathbf{1}\{D\geq d\}\mid X=x
\right)
}{
\mathrm{Var}(D\mid X=x)
}.
\]

\end{frame}


% ------------------------------------------------------------
\begin{frame}{Continuous Treatment + Conditional Ignorability}
\small

Suppose
\(\mathbb{E}[Y(d)\mid D,X]=\mathbb{E}[Y(d)\mid X]\).
Then \(m(d,x)=\mathbb{E}[Y(d)\mid X=x]\). Therefore,

\[
\boxed{
\beta^*
=
\frac{
\mathbb{E}
\left[
\mathrm{Var}(D\mid X)
\int
\omega(d\mid X)
\frac{\partial}{\partial d}
\mathbb{E}[Y(d)\mid X]
\,dd
\right]
}{
\mathbb{E}
[
\mathrm{Var}(D\mid X)
]
}.
}
\]

Again there are two layers of weighting:

\[
\boxed{
\underbrace{\mathrm{Var}(D\mid X)}_{\text{which covariate cells?}}
\times
\underbrace{\omega(d\mid X)}_{\text{which treatment margins?}}.
}
\]

\begin{block}{Interpretation}
A weighted average of conditional marginal causal effects.
\end{block}

\end{frame}


% ============================================================
% WHAT IS AND IS NOT CLEAN
% ============================================================

\begin{frame}{Where the Clean Results Stop}
\small

The previous covariate-adjusted formulas used the benchmark
\(D-\mathbb{E}[D\mid X]\).
Ordinary additive OLS instead uses
\(D-L(D\mid 1,X)\).

\medskip

If \(L(D\mid 1,X)\neq\mathbb{E}[D\mid X]\)
and treatment effects are heterogeneous, the coefficient generally
does \alert{not} reduce to the simple weighting formulas derived above.

\medskip

\begin{alertblock}{Status of the literature}
There are important characterizations for particular settings,
especially for binary treatment.

However, for an arbitrarily misspecified additive OLS regression
with heterogeneous effects, there is no single universal
interpretation such as

\[
\text{``OLS always equals a particular ATE.''}
\]

The interpretation depends on the specification and the source of
treatment variation.
\end{alertblock}

\scriptsize
See, e.g., S{\l}oczy\'nski (2022) and Lee and Han (2024).

\end{frame}


% ------------------------------------------------------------
\begin{frame}{Binary Treatment: Ordinary Additive OLS Can Reweight Effects}
\small

For binary treatment, modern work shows that when treatment effects
are heterogeneous, the coefficient from

\[
Y=\alpha+\beta D+X'\gamma+u
\]

can place implicit weights on different treatment-effect components
that differ substantially from their population shares.

\medskip

Therefore,

\[
\boxed{
\text{conditional ignorability alone}
\;\not\Rightarrow\;
\text{ordinary additive OLS coefficient}=ATE.
}
\]

The problem is not necessarily confounding.
The problem may instead be the estimand generated by the regression's
implicit weighting scheme.

\begin{block}{Modern interpretation}
Even after identification has been solved, one still needs to ask:
\[
\boxed{
\text{Which causal effects does this regression average?}
}
\]
\end{block}

\scriptsize
S{\l}oczy\'nski (2022), \emph{Review of Economics and Statistics}.

\end{frame}


% ============================================================
% MODERN TARGETS
% ============================================================

\begin{frame}{A More Modern Strategy: Choose the Estimand First}
\small

Rather than asking
``What causal parameter does this regression happen to estimate?'',
a modern approach is:

\[
\boxed{
\text{Scientific question}
\rightarrow
\text{causal estimand}
\rightarrow
\text{identification}
\rightarrow
\text{estimator}.
}
\]

Examples include:

\begin{itemize}
    \item Binary treatment:
    \[
    ATE,\quad ATT,\quad ATU,\quad
    \text{overlap-weighted ATE}.
    \]

    \item Multi-valued treatment:
    \(\mathbb{E}[Y(d_j)-Y(d_k)]\)
    or prespecified averages of treatment contrasts.

    \item Continuous treatment:
    \[
    \text{dose-response curve},
    \quad
    \text{shift effect},
    \quad
    ADE,
    \quad
    \text{weighted ADE}.
    \]
\end{itemize}

OLS may estimate one useful summary, but the regression coefficient
need not define the scientific target.

\end{frame}


% ------------------------------------------------------------
\begin{frame}{An Alternative Modern Object: Average Linear Regression Function}
\small

Graham and Pinto propose working with the conditional linear
projection within each covariate cell. Let

\[
b_0(x)
=
\frac{
\mathrm{Cov}(D,Y\mid X=x)
}{
\mathrm{Var}(D\mid X=x)
}.
\]

Instead of the pooled OLS weighting, consider

\[
\frac{
\mathbb{E}
[
\mathrm{Var}(D\mid X)b_0(X)
]
}{
\mathbb{E}
[
\mathrm{Var}(D\mid X)
]
}
\qquad\text{vs.}\qquad
\boxed{
\beta_{ALRF}
=
\mathbb{E}[b_0(X)].
}
\]

\begin{block}{Why this matters}
The choice between these two objects is itself an estimand choice.

Pooled regression automatically gives more weight to cells with
greater treatment variance.

The average linear regression function averages conditional
regression slopes over the covariate distribution instead.
\end{block}

\scriptsize
Graham and Pinto (2022), \emph{Journal of Econometrics}.

\end{frame}


% ------------------------------------------------------------
\begin{frame}{Continuous Treatments: The Frontier Is About the Target}
\small

For continuous treatment, the algebraic interpretation of OLS is
well understood:

\[
\boxed{
OLS
=
\text{a particular weighted average derivative}.
}
\]

What is \alert{not} uniquely determined is whether this is the
scientifically appropriate causal summary.

\medskip

Recent work studies alternative targets such as

\[
ADE
=
\mathbb{E}
\left[
\frac{\partial}{\partial d}
\mathbb{E}[Y(d)\mid X]
\bigg|_{d=D}
\right],
\]

as well as more general weighted average derivative effects.

\begin{alertblock}{Frontier status}
There is no single universally preferred scalar causal estimand for
a general continuous treatment.

The appropriate target depends on the intervention and the scientific
question.
\end{alertblock}

\scriptsize
See Graham and Pinto (2022) and Hines, Diaz-Ordaz, and Vansteelandt
(2026).

\end{frame}


% ============================================================
% SYNTHESIS
% ============================================================

\begin{frame}{Summary I: Without Covariates}
\small

\begin{columns}[T]
\begin{column}{0.32\textwidth}
\begin{block}{Binary \(D\)}
\[
\beta^{OLS}
=
\mathbb{E}[Y\mid 1]
-
\mathbb{E}[Y\mid 0].
\]

\medskip
Under exogeneity:
\[
\boxed{\beta^{OLS}=ATE.}
\]
\end{block}
\end{column}

\begin{column}{0.32\textwidth}
\begin{block}{Ordered discrete \(D\)}
\[
\beta^{OLS}
=
\sum_j w_j
\frac{
m(d_j)-m(d_{j-1})
}{
d_j-d_{j-1}
}.
\]

Under exogeneity: a weighted average of
\alert{causal treatment-margin slopes}.
\end{block}
\end{column}

\begin{column}{0.32\textwidth}
\begin{block}{Continuous \(D\)}
\[
\beta^{OLS}
=
\int \omega(d)m'(d)\,dd.
\]

\medskip
Under exogeneity: a weighted average of
\alert{marginal causal responses}.
\end{block}
\end{column}
\end{columns}

\bigskip

\[
\boxed{
\text{Same covariance ratio; the support of } D
\text{ determines what it averages.}
}
\]

\end{frame}


% ------------------------------------------------------------
\begin{frame}{Summary II: With Flexible Covariate Adjustment}
\small

The common formula is

\[
\boxed{
\beta^*
=
\frac{
\mathbb{E}
[
\mathrm{Var}(D\mid X)\beta(X)
]
}{
\mathbb{E}
[
\mathrm{Var}(D\mid X)
]
}.
}
\]

\begin{block}{Binary treatment}
\[
\beta(X)
=
\mathbb{E}[Y\mid 1,X]
-
\mathbb{E}[Y\mid 0,X].
\]

Under conditional ignorability:
\[
\boxed{
\beta^*
=
\text{overlap-weighted CATE}.
}
\]
\end{block}

\begin{block}{Ordered discrete treatment}
\[
\boxed{
\beta(X)
=
\text{weighted average across discrete treatment margins}.
}
\]
\end{block}

\begin{block}{Continuous treatment}
\[
\boxed{
\beta(X)
=
\text{weighted average of conditional derivatives}.
}
\]
\end{block}

\end{frame}


% ------------------------------------------------------------
\begin{frame}{The Main Lesson}
\small

The classical question is:
``What is the coefficient on \(D\)?''
The modern question is:

\[
\boxed{
\begin{array}{c}
\text{What population variation does the coefficient use,}\\
\text{and which heterogeneous effects does it average?}
\end{array}
}
\]

\medskip

OLS can be viewed as an \alert{averaging machine}:

\[
\boxed{
\text{OLS}
=
\text{selection of identifying variation}
+
\text{implicit weighting}.
}
\]

For binary treatment without covariates, this complexity is largely
hidden because there is only one treatment margin.
For multivalued or continuous treatment, the weighting problem
becomes explicit.

\bigskip

\[
\boxed{
\text{Whose effect?}
\qquad
\text{At which treatment margin?}
\qquad
\text{With what weight?}
}
\]

\end{frame}


% ------------------------------------------------------------
\begin{frame}{From ``Regression First'' to ``Estimand First''}
\small

A useful intellectual progression is

\[
\begin{array}{c}
\boxed{\text{Regression coefficient}}
\\[1mm]
\Downarrow
\\[1mm]
\boxed{\text{Population linear projection estimand}}
\\[1mm]
\Downarrow
\\[1mm]
\boxed{\text{Weighted average of conditional contrasts or slopes}}
\\[1mm]
\Downarrow\ \ \text{\scriptsize identification assumptions}
\\[1mm]
\boxed{\text{Weighted average of causal effects}}
\\[1mm]
\Downarrow
\\[1mm]
\boxed{\text{Ask whether this is the causal estimand we actually want.}}
\end{array}
\]

\medskip

This last step is the central modern shift:

\[
\boxed{
\text{Define the estimand first; choose the estimator second.}
}
\]

\end{frame}


% ============================================================
% REFERENCES
% ============================================================

\begin{frame}{Selected References}
\footnotesize

\begin{itemize}

    \item Graham, B. S., and C. C. de Xavier Pinto (2022).
    ``Semiparametrically Efficient Estimation of the Average Linear
    Regression Function.''
    \emph{Journal of Econometrics}, 226(1), 115--138.

    \item Hines, O. J., K. Diaz-Ordaz, and S. Vansteelandt (2026).
    ``Parameterizing the Effect of a Continuous Treatment Using
    Average Derivative Effects.''
    \emph{Biometrika}, 113(2).

    \item Ishimaru, S. (2024).
    ``Empirical Decomposition of the IV--OLS Gap with Heterogeneous
    and Nonlinear Effects.''
    \emph{Review of Economics and Statistics}, 106(2), 505--520.

    \item Lee, M.-J., and C. Han (2024).
    ``Ordinary Least Squares and Instrumental-Variables Estimators
    for Any Outcome and Heterogeneity.''
    \emph{Stata Journal}, 24(1), 72--92.

    \item S{\l}oczy\'nski, T. (2022).
    ``Interpreting OLS Estimands When Treatment Effects Are
    Heterogeneous: Smaller Groups Get Larger Weights.''
    \emph{Review of Economics and Statistics}, 104(3), 501--509.

    \item Yitzhaki, S. (1996).
    ``On Using Linear Regressions in Welfare Economics.''
    \emph{Journal of Business \& Economic Statistics}, 14(4), 478--486.

\end{itemize}

\end{frame}




\end{document}
